% ============================================================
%  第五讲
% ============================================================
\section{第五讲}

\subsection{Whittaker 系数}
\srcpage{15}

我们有如下箴言：

\begin{quote}
\(\Rep(\Gd)_{\mathrm{Ran}}\) 让我们把 \(\Bun_G\) 上一个层的所有 Whittaker 系数「粘」在一起。
\end{quote}

为解释这句话，先对一般 \(G\) 定义 Whittaker 系数。考虑
\[
\Bun^\Omega_N:=\Bun_B\times_{\Bun_T}\{\check\rho(\Omega)\},
\]
其中 \(\check\rho(\Omega):=2\check\rho(\Omega^{1/2})\)，\(\Omega^{1/2}\) 是线丛 \(\Omega\) 的一个固定平方根。
对每个正单根 \(\alpha_i\)，映射 \(N\to\Ga\) 诱导出映射
\[
\phi_i\colon\Bun^\Omega_N\to\Bun^\Omega_{\Ga}\to H^1(X,\Omega)\simeq\A^1.
\]
定义
\[
\phi:=\sum_{\alpha_i}\phi_i\colon\Bun^\Omega_N\to\A^1,
\]
并考虑对应
\[
\begin{tikzcd}
& \Bun^\Omega_N \arrow[dl,"p_N"'] \arrow[dr,"\phi"] & \\
\Bun_G & & \A^1
\end{tikzcd}
\]

与 \(\mathrm{PGL}_2\) 情形一样，我们作如下定义。

\begin{definition}\label{def:whittaker0}
设 \(\coeff_0\colon\Shv(\Bun_G)\to\Vect\) 为由下式给出的函子：
\[
\coeff_0(F):=\mathrm{C}^\bullet_c\bigl(\Bun^\Omega_N,\ p_N^!(F)\overset{!}{\otimes}\phi^!(\exp)\bigr).
\]
类似地，对每个 \(D=\sum\check\lambda_i x_i\)（其中 \(\lambda_i\in\check\Lambda^+\)），用
\(\check\rho(\Omega)(-D)\in\Bun_T\) 产生函子
\[
\coeff_D\colon\Shv(\Bun_G)\to\Vect.
\]
我们称它们为 \emph{Whittaker 系数函子}。
\yizhu{Whittaker 系数的经典原型是自守表示上的 Whittaker 函数：固定一个非平凡特征 \(\psi\) 于单幂根，
取矩阵系数的 \(\psi\)-分量。几何情形把「非平凡特征」换成 \(\exp\) 层（de Rham 的 \(e^t\) 或
\(\ell\)-进的 Artin--Schreier 层），把「积分」换成紧支集上同调。}{背景直觉·A：Frenkel--Gaitsgory--Vilonen, \emph{Whittaker patterns in the geometry of moduli spaces of bundles on curves}, Ann. of Math. 153 (2001)；B：nLab ``Whittaker function''}
\end{definition}

\begin{example}\label{ex:gm-whittaker}
对 \(G=\Gm\)，\(D\) 可以是任意除子，函子
\[
\coeff_D\colon\Shv(\Bun_{\Gm})\to\Vect
\]
就是取 \(\cO(-D)\) 处的 \(!\)-纤维。
\end{example}

\begin{example}\label{ex:pgl2-whittaker}
对 \(G=\mathrm{PGL}_2\)，这些函子回到前面讲次中用同样记号表示的函子。
\end{example}

下述结果称为几何 Casselman--Shalika 公式。

\begin{theorem}[Frenkel--Gaitsgory--Vilonen]\label{thm:casselman-shalika}
对 \(D=\sum\check\lambda_i x_i\)（\(\lambda_i\in\check\Lambda^+\)），有
\[
\coeff_D(F)\simeq\coeff_0\bigl(H_{\boxtimes V_{\check\lambda_i},x}(F)\bigr),
\]
其中回忆 \(H_{\boxtimes V_{\check\lambda_i},x}\) 是任意顺序复合 \(H_{V_{\check\lambda_i},x_i}\)。
\end{theorem}

注意有对象
\[
\boxtimes_i\bigl(V_{\check\lambda_i}\otimes\delta_{x_i}\bigr)
\in\Rep(\Gd)^{\otimes I}\otimes\Shv(X^I),
\]
它给出 \(\Rep(\Gd)_{\mathrm{Ran}}\) 中的一个对象。上述定理说的是：\(\coeff_D\) 等价于先作用该对象的 Hecke 作用，
再复合 \(\coeff_0\)。

\begin{definition}\label{def:normalized}
设 \(\cF_\sigma\) 是以 \(\sigma\) 为特征值的 Hecke 特征层。称 \(F\) 是\emph{归一化的}，
若它带有同构
\[
\coeff_0(F)\simeq\one.
\]
\end{definition}

设 \(\cF_\sigma\) 是归一化的 Hecke 特征层。由练习 \ref{ex:eigen-hecke}，有
\[
H_{\boxtimes V_{\check\lambda_i},x}(\cF_\sigma)
\simeq\cF_\sigma\otimes\bigotimes_i(V_{\check\lambda_i,\sigma})_{x_i}.
\]
由上述定理，这个条件蕴含
\[
\coeff_D(\cF_\sigma)
\simeq\coeff_0(\cF_\sigma)\otimes\bigotimes_i(V_{\check\lambda_i,\sigma})_{x_i}
\simeq\bigotimes_i(V_{\check\lambda_i,\sigma})_{x_i},
\]
其中最后一个同构是因为 \(\cF_\sigma\) 归一化。注意这推广了 \eqref{eq:3.1}。

现在考虑下述函子
\begin{equation}\label{eq:5.1}
\Rep(\Gd)_{\mathrm{Ran}}\otimes\Shv(\Bun_G)
\xrightarrow{\ \text{Hecke 作用}\ }\Shv(\Bun_G)
\xrightarrow{\ \coeff_0\ }\Vect.
\end{equation}
上述定理表明，这个函子可以看作把所有 Whittaker 系数装配成一个族。范畴
\(\Rep(\Gd)_{\mathrm{Ran}}\) 是自对偶的，配对函子为
\[
\Rep(\Gd)_{\mathrm{Ran}}\otimes\Rep(\Gd)_{\mathrm{Ran}}
\xrightarrow{\ -\star-\ }\Rep(\Gd)_{\mathrm{Ran}}\xrightarrow{\ \Gamma\ }\Vect,
\]
其中 \(-\star-\) 是 \(\Rep(\Gd)_{\mathrm{Ran}}\) 上的对称幺半结构，\(\Gamma:=\Hom(\mathrm{triv},-)\)。
于是函子 \eqref{eq:5.1} 可以改写成函子
\[
\coeff^{\mathrm{ult}}\colon\Shv(\Bun_G)\to\Rep(\Gd)_{\mathrm{Ran}},
\]
它由下述公式刻画：
\[
\Gamma\bigl(\coeff^{\mathrm{ult}}(F)\star G\bigr)
\simeq\coeff_0(G\cdot F),\qquad F\in\Shv(\Bun_G),\ G\in\Rep(\Gd)_{\mathrm{Ran}}.
\]
这里 \(G\cdot F\) 是 \(G\) 在 \(F\) 上的作用。

\begin{remark}\label{rem:q-expansion}
应当把 \(\coeff^{\mathrm{ult}}\) 看作自守层的「\(q\)-展开」的最佳版本。
\end{remark}

\begin{theorem}[Beraldo, Frenkel--Gaitsgory--Vilonon]\label{thm:coeff-ff}
对 \(G=\mathrm{GL}_n\) 或 \(\mathrm{PGL}_n\)，函子 \(\coeff^{\mathrm{ult}}\) 在
\(\Shv(\Bun_G)_{\mathrm{cusp}}\) 上是全忠实的。
\yizhu{定理 \ref{thm:coeff-ff} 的原文把最后一位作者拼成 ``Vilonon''，应为 ``Vilonen''。
这是排印错误，不影响数学内容。}{原文笔误·A：作者名}
\end{theorem}

对 \(\mathrm{PGL}_2\)，证明模仿 Fourier 反演；对 \(\mathrm{GL}_n\)，证明模仿
Piatetski-Shapiro 与 Shalika 的工作。

\subsection{范畴化的几何 Langlands 等价}
\srcpage{17}

现在我们在 de Rham 情形工作。考虑 \(\mathrm{LS}_{\Gd}=\mathrm{LS}^{\dR}_{\Gd}\)，
即 \(X\) 上（de Rham）\(\Gd\)-局部系统的模叠。存在典范对称幺半函子
\[
\Loc\colon\Rep(\Gd)_{\mathrm{Ran}}\to\QCoh(\mathrm{LS}_{\Gd}),
\]
它把 \(V\otimes\delta_x\in\Rep(\Gd)\otimes\DMod(X)\) 送到 \(\mathrm{ev}_x^*(V)\)。这里
\[
\mathrm{ev}_x\colon\mathrm{LS}_{\Gd}\to\mathrm{LS}_{\Gd}(D_x)\simeq B\Gd
\]
是赋值映射，并且我们把 \(\QCoh(B\Gd)\) 与 \(\Rep(\Gd)\) 等同。

\begin{exercise}\label{ex:ev-loc}
对 \(\sigma\in\mathrm{LS}_{\Gd}\)，把函子
\[
\mathrm{ev}_\sigma\colon\Rep(\Gd)_{\mathrm{Ran}}\to\Vect
\]
等同于复合
\[
\Rep(\Gd)_{\mathrm{Ran}}\xrightarrow{\ \Loc\ }\QCoh(\mathrm{LS}_{\Gd})
\xrightarrow{\ (-)^!_\sigma\ }\Vect.
\]
\end{exercise}

\begin{theorem}[Lurie, Gaitsgory--Rozenblyum]\label{thm:loc-adjoint}
函子 \(\Loc\) 有全忠实的右伴随
\[
\QCoh(\mathrm{LS}_{\Gd})\hookrightarrow\Rep(\Gd)_{\mathrm{Ran}}.
\]
\end{theorem}

\begin{theorem}[Drinfeld--Gaitsgory]\label{thm:action-factors}
\(\Rep(\Gd)_{\mathrm{Ran}}\) 在 \(\DMod(\Bun_G)\) 上的作用经 \(\Loc\) 分解。
\end{theorem}

\begin{corollary}\label{cor:coeff-factors}
函子 \(\coeff^{\mathrm{ult}}\) 分解为
\[
\DMod(\Bun_G)\xrightarrow{\ \LG_{\mathrm{coarse}}\ }\QCoh(\mathrm{LS}_{\Gd})
\hookrightarrow\Rep(\Gd)_{\mathrm{Ran}}.
\]
\end{corollary}

范畴化几何 Langlands 的想法是：函子 \(\LG_{\mathrm{coarse}}\) 「几乎」是等价。注意它不可能是等价，
因为对非交换的 \(G\)，它把常层送到 \(0\)。

\begin{theorem}[范畴化几何 Langlands，第一版]\label{thm:categorical-1}
函子 \(\LG_{\mathrm{coarse}}\) 诱导等价
\[
\DMod(\Bun_G)_{\mathrm{cusp}}\simeq\QCoh(\mathrm{LS}^{\mathrm{irred}}_{\Gd}),
\]
其中 \(\mathrm{LS}^{\mathrm{irred}}_{\Gd}\subset\mathrm{LS}_{\Gd}\) 是不可约 \(\Gd\)-局部系统的轨迹。
\end{theorem}

\begin{corollary}\label{cor:coeff-ff-cusp}
\(\coeff^{\mathrm{ult}}\) 在尖点子范畴上是全忠实的。
\end{corollary}

Arinkin--Gaitsgory 对 \(\LG_{\mathrm{coarse}}\) 还有一个修正。回忆：对足够好的叠 \(Y\)，有
\[
\QCoh(Y)\simeq\Ind(\Perf(Y)),
\]
其中 \(\Perf(Y)\) 是局部有界、有限秩投射模复形构成的全子范畴。我们还可以考虑
\[
\IndCoh(Y):=\Ind(\Coh(Y)),
\]
其中 \(\Coh(Y)\) 是局部有界、有限生成模复形构成的范畴。有包含
\[
\Perf(Y)\subset\Coh(Y),
\]
当 \(Y\) 有奇点时，这个包含是严格的。

现在可以定义子范畴
\[
\Coh_{\mathrm{Nilp}}(\mathrm{LS}_{\Gd})\subset\Coh(\mathrm{LS}_{\Gd}),
\]
它由所有抛物子群 \(\check P\) 及其 Levi 商 \(\check M\) 的下述像生成：
\[
\Perf(\mathrm{LS}_{\check M})\xrightarrow{\ q^* \ }\Perf(\mathrm{LS}_{\check P})
\xrightarrow{\ p_* \ }\Coh(\mathrm{LS}_{\Gd}).
\]
这里 \(q\colon\mathrm{LS}_{\check P}\to\mathrm{LS}_{\check M}\) 是有限 Tor 幅度的映射，
\(p\colon\mathrm{LS}_{\check P}\to\mathrm{LS}_{\Gd}\) 是固有映射，因此上述函子有定义。
\yizhu{\(\IndCoh_{\mathrm{Nilp}}\) 的下标 ``Nilp'' 指幂零奇异支撑。Arinkin--Gaitsgory 的修正动机是：
\(\Bun_G\) 一般有奇点，\(\QCoh\) 不是正确的目标；\(\IndCoh\) 加上幂零奇异支撑条件才是。
Geometric Langlands 的最终陈述因此写成 \(\DMod(\Bun_G)\simeq\IndCoh_{\mathrm{Nilp}}(\mathrm{LS}_{\Gd})\)，
而不是与 \(\QCoh\) 的等价。}{背景直觉·A：Arinkin--Gaitsgory, \emph{Singular support of coherent sheaves}; B：nLab ``geometric Langlands correspondence''}

\begin{remark}\label{rem:cohnilp}
子范畴 \(\Coh_{\mathrm{Nilp}}(\mathrm{LS}_{\Gd})\subset\Coh(\mathrm{LS}_{\Gd})\)
可以用凝聚层的奇异支撑理论内蕴地定义。
\end{remark}

\begin{theorem}\label{thm:LG}
存在唯一函子
\[
\LG\colon\DMod(\Bun_G)\to\IndCoh_{\mathrm{Nilp}}(\mathrm{LS}_{\Gd}),
\]
它满足一个技术条件\yizhu{定理 \ref{thm:LG} 的技术条件是：把源范畴中的紧对象送到目标 \(t\)-结构中有下界的对象。
这是原文脚注 6 的内容；译文中置于脚注，正文不重复。}{记号·A：原文脚注 6}，并使得下图交换：
\[
\begin{tikzcd}
\DMod(\Bun_G) \arrow[r,"\LG"] \arrow[dr,"\LG_{\mathrm{coarse}}"'] & \IndCoh_{\mathrm{Nilp}}(\mathrm{LS}_{\Gd}) \arrow[d,"\Psi"] \\
& \QCoh(\mathrm{LS}_{\Gd}),
\end{tikzcd}
\]
其中 \(\Psi\) 是包含 \(\Coh_{\mathrm{Nilp}}\subset\Coh\subset\QCoh\) 的 ind-延拓。
\end{theorem}

\begin{theorem}[范畴化几何 Langlands，最终版]\label{thm:categorical-ultimate}
函子 \(\LG\) 是等价。
\end{theorem}

\begin{remark}\label{rem:proof}
证明用到 de Rham 情形的特殊性，例如用临界水平的 Kac--Moody 局部化来证明关于
\(\DMod(\Bun_G)\) 的「几何」陈述。更严肃地说，我们用到 de Rham 模叠
\(\mathrm{LS}^{\dR}_{\Gd}\) 的全局（代数）函数很少这一事实。
尽管如此，我们仍可以通过某种意义上的「Riemann--Hilbert」得到 Betti 版本
（以及 \(\operatorname{char}(k)=0\) 时的 \(\ell\)-进版本）。
\end{remark}

\begin{remark}\label{rem:concrete}
上一讲的「具体」定理 \ref{thm:gl-general} 实际上用到了全部范畴断言。
\end{remark}
